\section{The Theorems of Fermat, Euler and Wilson}

\begin{theorem}[Euler]\label{thm:euler_theorem}
    Let $m \ge 1$ and $(a,m) = 1$. Then
    \begin{equation*}
        a^{\phi(m)} \equiv 1 \pmod m.
    \end{equation*}
\end{theorem}
\begin{proof}
    Let $r_1, \dots, r_{\phi(m)}$ be a reduced residue system modulo $m$. Each $ar_i$ is coprime to $m$, and $ar_i \equiv ar_j$ implies $r_i \equiv r_j$ by \cref{cor:cancellation}. Hence $ar_1, \dots, ar_{\phi(m)}$ is again a reduced residue system, and
    \begin{equation*}
        \prod_{i=1}^{\phi(m)} (ar_i) \equiv \prod_{i=1}^{\phi(m)} r_i \pmod m.
    \end{equation*}
    The left-hand side is $a^{\phi(m)} \prod_i r_i$, and $\prod_i r_i$ is coprime to $m$, so we can cancel it.
\end{proof}

\begin{corollary}[Fermat's little theorem]\label{cor:fermat}
    Let $p$ be prime. If $p \nmid a$, then $a^{p-1} \equiv 1 \pmod p$. For every $a \in \mathbb{Z}$, $a^p \equiv a \pmod p$.
\end{corollary}
\begin{proof}
    The first statement is \cref{thm:euler_theorem} with $\phi(p) = p - 1$. Multiplying by $a$ gives the second when $p \nmid a$, and it is trivial when $p \mid a$.
\end{proof}

\begin{theorem}[Wilson]\label{thm:wilson}
    An integer $n \ge 2$ is prime if and only if $(n-1)! \equiv -1 \pmod n$.
\end{theorem}
\begin{proof}
    Let $p$ be prime. The case $p = 2$ is clear, so assume $p$ odd. If $x^2 \equiv 1 \pmod p$, then $p \mid (x-1)(x+1)$, so $x \equiv \pm 1$ by \cref{cor:euclid_lemma}. Thus $1$ and $p - 1$ are the only elements of $\{1, \dots, p-1\}$ equal to their own inverse modulo $p$. The remaining $p - 3$ elements split into pairs $\{x, x^{-1}\}$ with product $1$, so
    \begin{equation*}
        (p-1)! \equiv 1 \cdot (p-1) \equiv -1 \pmod p.
    \end{equation*}

    Conversely, let $n$ be composite. Then $n$ has a divisor $d$ with $1 < d < n$, so $d \mid (n-1)!$. If $(n-1)! \equiv -1 \pmod n$, then $d$ would also divide $(n-1)! + 1$, hence divide $1$, which is impossible.
\end{proof}
